

\documentclass[11pt,a4paper]{article}
\usepackage[utf8]{inputenc}
\usepackage[T1]{fontenc}
\usepackage[french]{babel}
\usepackage{lmodern}
\usepackage{microtype}
\usepackage[a4paper,margin=1.65cm,headheight=15pt]{geometry}
\usepackage{amsmath,amssymb,mathtools}
\usepackage{array,tabularx,multicol,booktabs}
\usepackage{enumitem}
\usepackage[most]{tcolorbox}
\usepackage{xcolor}
\usepackage{tikz}
\usetikzlibrary{arrows.meta,calc,positioning}
\usepackage{pdflscape}
\usepackage{fancyhdr}
\usepackage{hyperref}
\usepackage{pagecolor}
\usepackage{needspace}

\definecolor{fondpage}{RGB}{247,248,250}
\definecolor{encre}{RGB}{31,35,40}
\definecolor{bleu}{RGB}{40,91,135}
\definecolor{vert}{RGB}{55,125,92}
\definecolor{orange}{RGB}{178,105,42}
\definecolor{violet}{RGB}{101,77,145}
\definecolor{rouge}{RGB}{170,72,72}
\definecolor{gris}{RGB}{86,91,99}
\definecolor{bleufonce}{RGB}{28,55,120}
\definecolor{bleuclair}{RGB}{238,243,255}
\definecolor{grisclair}{RGB}{245,245,245}

\hypersetup{colorlinks=true,linkcolor=bleufonce,urlcolor=bleufonce,
pdftitle={Parcours de revision -- Fonction exponentielle -- Premiere specialite},pdfauthor={}}

\setlength{\parindent}{0pt}
\setlength{\parskip}{0.25em}
\renewcommand{\arraystretch}{1.13}
\setlength{\tabcolsep}{4pt}
\setlist[itemize]{leftmargin=1.25em,itemsep=0.15em,topsep=0.2em}
\setlist[enumerate,1]{label=\textbf{\arabic*.}, leftmargin=1.05cm, itemsep=0.35em, topsep=0.2em}
\setlist[enumerate,2]{label=\textbf{\alph*.}, leftmargin=0.85cm, itemsep=0.25em, topsep=0.15em}

\pagestyle{fancy}
\fancyhf{}
\cfoot{\thepage}

\newcommand{\R}{\mathbb{R}}
\newcommand{\N}{\mathbb{N}}
\newcommand{\sourceq}[1]{ {\scriptsize\textcolor{bleufonce}{(site Premiere q#1)}}}
\newcounter{question}
\newcounter{reponse}

\newenvironment{blocquestions}[1]{%
  \begin{tcolorbox}[enhanced,breakable,colback=bleuclair,colframe=bleufonce,boxrule=0.6pt,arc=2mm,left=2mm,right=2mm,top=2mm,bottom=1.5mm,title={\bfseries #1},coltitle=white,colbacktitle=bleufonce,before skip=3mm,after skip=3mm, fontupper=\large]
  \large
  \begin{enumerate}[leftmargin=*,label=\textbf{\arabic*.},itemsep=0.82em,topsep=0.5em]
  \setcounter{enumi}{\value{question}}
}{%
  \setcounter{question}{\value{enumi}}
  \end{enumerate}
  \end{tcolorbox}
}

\newenvironment{blocreponses}[1]{%
  \begin{tcolorbox}[enhanced,breakable,colback=bleuclair,colframe=bleufonce,boxrule=0.6pt,arc=2mm,left=2mm,right=2mm,top=2mm,bottom=1.5mm,title={\bfseries #1},coltitle=white,colbacktitle=bleufonce,before skip=3mm,after skip=3mm, fontupper=\large]
  \large
  \begin{enumerate}[leftmargin=*,label=\textbf{\arabic*.},itemsep=0.42em,topsep=0.5em]
  \setcounter{enumi}{\value{reponse}}
}{%
  \setcounter{reponse}{\value{enumi}}
  \end{enumerate}
  \end{tcolorbox}
}

\newtcolorbox{correctionbox}[1]{enhanced,breakable,colback=grisclair,colframe=black!55,boxrule=0.55pt,arc=2mm,left=2mm,right=2mm,top=2mm,bottom=1.5mm,title=\textbf{#1},coltitle=black,colbacktitle=black!12,before skip=2mm,after skip=2mm}
\newtcolorbox{methodebox}[1]{enhanced,breakable,colback=white,colframe=bleu!70!black,boxrule=0.55pt,arc=2mm,left=2mm,right=2mm,top=2mm,bottom=1.5mm,title=\textbf{#1},coltitle=white,colbacktitle=bleu!70!black,before skip=2mm,after skip=2mm}
\newtcolorbox{exoBox}[2]{enhanced,breakable,colback=white,colframe=#1!70!black,boxrule=0.55pt,arc=2mm,left=2mm,right=2mm,top=2mm,bottom=1.5mm,title=\textbf{#2},coltitle=white,colbacktitle=#1!70!black,before skip=3mm,after skip=3mm}

\newcommand{\titreSujet}[2]{%
  \subsection*{#1}
  \addcontentsline{toc}{subsection}{#1}
  \textit{Référence / inspiration : #2}\par\medskip\hrule\medskip
}
\newcommand{\qcm}[4]{%
\begin{center}\begin{tabularx}{0.96\linewidth}{@{}XX@{}}
\textbf{a.} #1 & \textbf{b.} #2\\[1mm]
\textbf{c.} #3 & \textbf{d.} #4
\end{tabularx}\end{center}}

\tcbset{enhanced,arc=2mm,outer arc=2mm,boxrule=0.42pt,left=1.15mm,right=1.15mm,top=0.75mm,bottom=0.75mm,boxsep=0.35mm,before skip=0.8mm,after skip=0.8mm,fonttitle=\bfseries\footnotesize,coltitle=encre,before upper=\raggedright\fontsize{7.15}{7.65}\selectfont}
\newtcolorbox{fichebox}[3][]{colback=white,colframe=#2!72!black,colbacktitle=#2!18,coltitle=#2!50!black,borderline west={0.9mm}{0pt}{#2!78!black},title={#3},drop fuzzy shadow=black!5,#1}
\newcommand{\tagcol}[2]{\begin{tcolorbox}[enhanced,colback=#1!11,colframe=#1!45!black,boxrule=0.42pt,arc=2mm,left=1mm,right=1mm,top=0.45mm,bottom=0.45mm,before skip=0mm,after skip=0.85mm,fontupper=\bfseries\scriptsize,colupper=#1!55!black]\centering #2\end{tcolorbox}}
\newtcolorbox{panel}[1]{colback=fondpage,colframe=fondpage,boxrule=0pt,arc=0mm,left=0mm,right=0mm,top=0mm,bottom=0mm,height=168mm,valign=top,before skip=0mm,after skip=0mm}

\newcommand{\titrepage}{\begin{tcolorbox}[colback=white,colframe=bleu!70!black,boxrule=0.65pt,arc=3mm,left=2.5mm,right=2.5mm,top=1.15mm,bottom=1.15mm,before skip=0mm,after skip=1.4mm,drop fuzzy shadow=black!10]
{\Large\bfseries Parcours de révision -- Fonction exponentielle}\hfill {\normalsize Première spécialité -- sans calculatrice}
\end{tcolorbox}}
\newcommand{\courbeExp}{%
\begin{center}
\begin{tikzpicture}[scale=0.52]
\draw[->,gray] (-2.1,0)--(2.8,0) node[right]{$x$};
\draw[->,gray] (0,-0.2)--(0,4.4) node[above]{$y$};
\foreach \x in {-2,-1,0,1,2}{\draw[gray] (\x,0.06)--(\x,-0.06) node[below] {\scriptsize $\x$};}
\foreach \y/\lab in {1/1,2/2,3/3}{\draw[gray] (0.06,\y)--(-0.06,\y) node[left] {\scriptsize $\lab$};}
\draw[very thick,bleu,domain=-2:1.45,samples=90] plot (\x,{exp(\x)});
\node[bleu] at (1.65,3.5) {$y=e^x$};
\draw[dashed,gray] (1,0)--(1,2.718);
\fill (1,2.718) circle (1.5pt) node[right] {\scriptsize $e$};
\end{tikzpicture}
\end{center}}

\newcommand{\tableVarSimple}[5]{%
\begin{center}
\begin{tikzpicture}[x=1cm,y=0.7cm]
\draw[fill=grisclair] (0,2) rectangle (8,3);
\draw (0,0) rectangle (8,3);
\draw (1.45,0)--(1.45,3);
\draw (0,1)--(8,1); \draw (0,2)--(8,2);
\node at (0.72,2.5) {$x$};
\node at (0.72,1.5) {$f'(x)$};
\node at (0.72,0.5) {$f$};
\node at (1.85,2.5) {$#1$};
\node at (4.65,2.5) {$#2$};
\node at (7.45,2.5) {$#3$};
\node at (2.9,1.5) {$#4$};
\node at (4.65,1.5) {$0$};
\node at (6.4,1.5) {$#5$};
\end{tikzpicture}
\end{center}}

\newcommand{\ficheRecap}{%
\begin{landscape}
\pagecolor{fondpage}
\thispagestyle{empty}
\titrepage
\begin{tabularx}{\linewidth}{@{}X@{\hspace{1.8mm}}X@{\hspace{1.8mm}}X@{\hspace{1.8mm}}X@{}}
\begin{panel}{}
\tagcol{bleu}{Définition et courbe}
\begin{fichebox}{bleu}{Définition}
La fonction exponentielle est l'unique fonction $\exp$ dérivable sur $\mathbb R$ telle que :
\[
\exp(0)=1 \quad\text{et}\quad \exp'=\exp.
\]
On note :
\[
\exp(x)=e^x.
\]
\end{fichebox}
\begin{fichebox}{bleu}{Valeurs de référence}
\[
e^0=1,\qquad e^1=e.
\]
Le nombre $e$ est environ égal à $2{,}718$.
\end{fichebox}
\begin{fichebox}{bleu}{Courbe et signe}
\courbeExp
Pour tout réel $x$ :
\[
e^x>0.
\]
\end{fichebox}
\end{panel}
&
\begin{panel}{}
\tagcol{vert}{Propriétés algébriques}
\begin{fichebox}{vert}{Formules}
Pour tous réels $a$ et $b$ :
\[
e^{a+b}=e^a e^b,
\qquad
 e^{a-b}=\frac{e^a}{e^b}.
\]
\[
(e^a)^n=e^{na},
\qquad
 e^{-a}=\frac1{e^a}.
\]
\end{fichebox}
\begin{fichebox}{vert}{Attention}
On ne développe pas $e^{a+b}$ :
\[
e^{a+b}\neq e^a+e^b.
\]
\end{fichebox}
\end{panel}
&
\begin{panel}{}
\tagcol{orange}{Dériver et varier}
\begin{fichebox}{orange}{Dérivées de référence}
\[
(e^x)'=e^x.
\]
Pour tous réels $a$ et $b$ :
\[
\bigl(e^{ax+b}\bigr)'=a e^{ax+b}.
\]
\end{fichebox}
\begin{fichebox}{orange}{Produit avec exponentielle}
Pour dériver un produit, on utilise :
\[
(uv)'=u'v+uv'.
\]
Exemple :
\[
\bigl(xe^x\bigr)'=e^x+xe^x=(x+1)e^x.
\]
\end{fichebox}
\end{panel}
&
\begin{panel}{}
\tagcol{violet}{Équations et inéquations}
\begin{fichebox}{violet}{Croissance}
La fonction exponentielle est strictement croissante sur $\mathbb R$.\\[0.3em]
Donc :
\[
e^a=e^b \Longleftrightarrow a=b.
\]
\[
e^a<e^b \Longleftrightarrow a<b.
\]
\end{fichebox}
\begin{fichebox}{violet}{Méthode}
Pour résoudre une équation du type :
\[
e^{2x-1}=e^{x+3},
\]
on compare les exposants :
\[
2x-1=x+3.
\]
\end{fichebox}
\end{panel}
\end{tabularx}
\clearpage
\nopagecolor
\end{landscape}}

\begin{document}
\begin{center}
{\LARGE\bfseries Corrigé -- Parcours de révision -- Fonction exponentielle}\\[1mm]
{\large Première générale -- spécialité mathématiques}\\[1mm]
{\normalsize Sans calculatrice}
\end{center}
\tableofcontents\newpage
\phantomsection\addcontentsline{toc}{section}{Fiche récapitulative}
\ficheRecap

\section{Corrigé des questions flash}
\setcounter{reponse}{0}
\begin{blocreponses}{Réponses 1 à 12}
\item Calculer $e^0$.\quad \textbf{Réponse :} $1$.
\item Calculer $e^1$.\quad \textbf{Réponse :} $e$.
\item Simplifier $e^{x+2}e^{3}$.\quad \textbf{Réponse :} $e^{x+5}$.
\item Simplifier $\dfrac{e^{x+4}}{e^x}$.\quad \textbf{Réponse :} $e^4$.
\item Simplifier $e^{2x}e^{1-x}$.\quad \textbf{Réponse :} $e^{x+1}$.
\item Simplifier $\dfrac{e^{3x}}{e^x}$.\quad \textbf{Réponse :} $e^{2x}$.
\item Simplifier $(e^x)^4$.\quad \textbf{Réponse :} $e^{4x}$.
\item Simplifier $e^{-x}e^{2x}$.\quad \textbf{Réponse :} $e^x$.
\item Dériver $f(x)=e^x$.\quad \textbf{Réponse :} $f'(x)=e^x$.
\item Dériver $f(x)=3e^x$.\quad \textbf{Réponse :} $f'(x)=3e^x$.
\item Dériver $f(x)=e^{2x}$.\quad \textbf{Réponse :} $f'(x)=2e^{2x}$.
\item Dériver $f(x)=e^{-x}$.\quad \textbf{Réponse :} $f'(x)=-e^{-x}$.
\end{blocreponses}

\begin{blocreponses}{Réponses 13 à 24}
\item Dériver $f(x)=e^{3x-1}$.\quad \textbf{Réponse :} $f'(x)=3e^{3x-1}$.
\item Dériver $f(x)=2e^{x}-x$.\quad \textbf{Réponse :} $f'(x)=2e^x-1$.
\item Dériver $f(x)=xe^x$.\quad \textbf{Réponse :} $f'(x)=(x+1)e^x$.
\item Dériver $f(x)=(x-2)e^x$.\quad \textbf{Réponse :} $f'(x)=(x-1)e^x$.
\item Résoudre $e^x=1$.\quad \textbf{Réponse :} $x=0$.
\item Résoudre $e^{x+2}=e^5$.\quad \textbf{Réponse :} $x=3$.
\item Résoudre $e^{2x}=e^{x+4}$.\quad \textbf{Réponse :} $x=4$.
\item Résoudre $e^{3x-1}=e^{x+5}$.\quad \textbf{Réponse :} $x=3$.
\item Résoudre $e^{x}<e^4$.\quad \textbf{Réponse :} $x<4$.
\item Résoudre $e^{2x-1}\ge e^3$.\quad \textbf{Réponse :} $x\ge2$.
\item Donner le signe de $e^x$.\quad \textbf{Réponse :} Strictement positif.
\item Donner le sens de variation de $e^x$.\quad \textbf{Réponse :} Strictement croissante.
\end{blocreponses}

\begin{blocreponses}{Réponses 25 à 36}
\item Donner le signe de $(x-3)e^x$.\quad \textbf{Réponse :} Signe de $x-3$.
\item Donner le signe de $-(x+1)e^x$.\quad \textbf{Réponse :} Signe de $-(x+1)$.
\item Factoriser $e^{2x}-e^x$.\quad \textbf{Réponse :} $e^x(e^x-1)$.
\item Résoudre $e^x(e^x-1)=0$.\quad \textbf{Réponse :} $x=0$.
\item Dériver $f(x)=x+e^{1-x}$.\quad \textbf{Réponse :} $f'(x)=1-e^{1-x}$.
\item Dériver $f(x)=(2x+1)e^x$.\quad \textbf{Réponse :} $f'(x)=(2x+3)e^x$.
\item Si $f'(x)=(x-1)e^x$, donner le signe de $f'(x)$.\quad \textbf{Réponse :} Signe de $x-1$.
\item Si $f'(x)=-2e^{x}$, donner le sens de variation de $f$.\quad \textbf{Réponse :} $f$ est décroissante.
\item Calculer $e^2/e^{-1}$.\quad \textbf{Réponse :} $e^3$.
\item Simplifier $e^{x-2}\times e^{2-x}$.\quad \textbf{Réponse :} $e^0=1$.
\item Résoudre $e^{4-x}=e^{1+x}$.\quad \textbf{Réponse :} $x=\frac32$.
\item Résoudre $e^{x+1}\le e^{2x-3}$.\quad \textbf{Réponse :} $x\ge4$.
\end{blocreponses}


\section{Corrigé du QCM}
\begin{correctionbox}{Réponses du QCM}
\begin{center}
\renewcommand{\arraystretch}{1.25}
\begin{tabular}{c|cccccccc}
Question & 1&2&3&4&5&6&7&8\\\hline
Réponse & b&b&b&c&c&a&a&c\\
\end{tabular}
\end{center}
\end{correctionbox}

\section{Corrigé des sujets type bac}


\titreSujet{Sujet 1 -- Corrigé}{corrigé}
\begin{enumerate}
\item On calcule :
\[
f(x)-g(x)=e^{2x}+e^2-(1+e^2)e^x.
\]
Donc :
\[
f(x)-g(x)=e^{2x}-(1+e^2)e^x+e^2.
\]

\item On cherche à factoriser :
\[
X^2-(1+e^2)X+e^2.
\]
Or :
\[
(X-1)(X-e^2)=X^2-(1+e^2)X+e^2.
\]
Ainsi :
\[
X^2-(1+e^2)X+e^2=(X-1)(X-e^2).
\]

\item En posant $X=e^x$, puis en utilisant la factorisation précédente, on obtient :
\[
f(x)-g(x)=(e^x-1)(e^x-e^2).
\]

\item Les points d'intersection vérifient $f(x)=g(x)$, donc :
\[
f(x)-g(x)=0.
\]
Ainsi :
\[
(e^x-1)(e^x-e^2)=0.
\]
Donc $e^x=1=e^0$ ou $e^x=e^2$. Par conséquent :
\[
x=0\quad\text{ou}\quad x=2.
\]

\item Pour $x=0$, on a :
\[
f(0)=e^0+e^2=1+e^2.
\]
Le premier point d'intersection est donc :
\[
A(0;1+e^2).
\]
Pour $x=2$, on a :
\[
f(2)=e^4+e^2.
\]
Le second point d'intersection est donc :
\[
B(2;e^4+e^2).
\]

\item Comme la fonction exponentielle est strictement croissante :
\[
e^x-1<0 \Longleftrightarrow x<0
\qquad\text{et}\qquad
e^x-e^2<0 \Longleftrightarrow x<2.
\]
On en déduit le tableau de signe suivant :
\begin{center}
\begin{tikzpicture}[x=1cm,y=0.62cm]
\draw[fill=grisclair] (0,3) rectangle (9.2,4);
\draw (0,0) rectangle (9.2,4);
\draw (1.7,0)--(1.7,4);
\draw (0,1)--(9.2,1); \draw (0,2)--(9.2,2); \draw (0,3)--(9.2,3);
\node at (0.85,3.5) {$x$};
\node at (0.85,2.5) {$e^x-1$};
\node at (0.85,1.5) {$e^x-e^2$};
\node at (0.85,0.5) {$f(x)-g(x)$};
\node at (2.2,3.5) {$-\infty$}; \node at (4.5,3.5) {$0$}; \node at (6.8,3.5) {$2$}; \node at (8.7,3.5) {$+\infty$};
\node at (3.2,2.5) {$-$}; \node at (4.5,2.5) {$0$}; \node at (5.65,2.5) {$+$}; \node at (7.75,2.5) {$+$};
\node at (3.2,1.5) {$-$}; \node at (5.65,1.5) {$-$}; \node at (6.8,1.5) {$0$}; \node at (7.75,1.5) {$+$};
\node at (3.2,0.5) {$+$}; \node at (4.5,0.5) {$0$}; \node at (5.65,0.5) {$-$}; \node at (6.8,0.5) {$0$}; \node at (7.75,0.5) {$+$};
\end{tikzpicture}
\end{center}

\item Comme $f(x)-g(x)$ représente l'écart entre les ordonnées des deux courbes :
\begin{itemize}
\item $\mathcal C_f$ est au-dessus de $\mathcal C_g$ sur $]-\infty;0[$ et sur $]2;+\infty[$ ;
\item les deux courbes se coupent pour $x=0$ et $x=2$ ;
\item $\mathcal C_f$ est en dessous de $\mathcal C_g$ sur $]0;2[$.
\end{itemize}

\item D'après le tableau de signe :
\[
f(x)\le g(x) \Longleftrightarrow f(x)-g(x)\le 0.
\]
Donc :
\[
\boxed{x\in[0;2]}.
\]

\item On remarque que :
\[
f(1)=e^2+e^2=2e^2
\]
et :
\[
g(1)=(1+e^2)e=e+e^3.
\]
Or $1\in[0;2]$, donc $f(1)\le g(1)$. Ainsi :
\[
\boxed{2e^2\le e+e^3}.
\]
Comme $1$ n'est pas une abscisse d'intersection, l'inégalité est même stricte :
\[
\boxed{2e^2<e+e^3}.
\]
\end{enumerate}

\titreSujet{Sujet 2 -- Corrigé}{corrigé}
\begin{enumerate}
\item On calcule :
\[
f(0)=(0-2)e^0=-2.
\]
\item On dérive un produit :
\[
f'(x)=1\times e^x+(x-2)e^x=(x-1)e^x.
\]
\item Comme $e^x>0$, le signe de $f'(x)$ est celui de $x-1$.
\[
f'(x)<0 \text{ sur }[-2;1[,\quad f'(1)=0,\quad f'(x)>0 \text{ sur }]1;4].
\]
\item Tableau de variations :
\begin{center}
\begin{tikzpicture}[x=1cm,y=0.72cm]
\draw[fill=grisclair] (0,2) rectangle (8.5,3);
\draw (0,0) rectangle (8.5,3);
\draw (1.4,0)--(1.4,3);
\draw (0,1)--(8.5,1); \draw (0,2)--(8.5,2);
\node at (0.7,2.5) {$x$}; \node at (0.7,1.5) {$f'(x)$}; \node at (0.7,0.5) {$f$};
\node at (1.8,2.5) {$-2$}; \node at (4.95,2.5) {$1$}; \node at (8.0,2.5) {$4$};
\node at (3.2,1.5) {$-$}; \node at (4.95,1.5) {$0$}; \node at (6.55,1.5) {$+$};
\node at (1.95,0.78) {$-4e^{-2}$}; \node at (5.0,0.18) {$-e$}; \node at (7.8,0.78) {$2e^4$};
\draw[-{Latex[length=2mm]},thick] (2.25,0.72)--(4.55,0.25);
\draw[-{Latex[length=2mm]},thick] (5.35,0.25)--(7.35,0.72);
\end{tikzpicture}
\end{center}
\item Le minimum est atteint pour $x=1$ et vaut :
\[
f(1)=(1-2)e=-e.
\]
\item Comme $e^x>0$, le signe de $f(x)$ est celui de $x-2$.\
Donc $f(x)<0$ sur $[-2;2[$, $f(2)=0$, et $f(x)>0$ sur $]2;4]$.
\item La tangente au point d'abscisse $0$ a pour équation :
\[
y=f'(0)(x-0)+f(0).
\]
Or $f(0)=-2$ et $f'(0)=(0-1)e^0=-1$.\
Donc :
\[
y=-x-2.
\]
\end{enumerate}


\titreSujet{Sujet 3 -- Corrigé}{corrigé}
\begin{enumerate}
\item On calcule :
\[
C(0)=0\times e^1=0,
\qquad
C(1)=1\times e^0=1.
\]
\item On dérive le produit $t e^{1-t}$ :
\[
C'(t)=1\times e^{1-t}+t\times(-e^{1-t}).
\]
Donc :
\[
C'(t)=(1-t)e^{1-t}.
\]
\item Comme $e^{1-t}>0$, le signe de $C'(t)$ est celui de $1-t$.\
Donc $C'(t)>0$ sur $[0;1[$, $C'(1)=0$, et $C'(t)<0$ sur $]1;5]$.
\item Tableau de variations :
\begin{center}
\begin{tikzpicture}[x=1cm,y=0.72cm]
\draw[fill=grisclair] (0,2) rectangle (8.5,3);
\draw (0,0) rectangle (8.5,3);
\draw (1.4,0)--(1.4,3);
\draw (0,1)--(8.5,1); \draw (0,2)--(8.5,2);
\node at (0.7,2.5) {$t$}; \node at (0.7,1.5) {$C'(t)$}; \node at (0.7,0.5) {$C$};
\node at (1.8,2.5) {$0$}; \node at (4.95,2.5) {$1$}; \node at (8.0,2.5) {$5$};
\node at (3.2,1.5) {$+$}; \node at (4.95,1.5) {$0$}; \node at (6.55,1.5) {$-$};
\node at (1.9,0.18) {$0$}; \node at (5.0,0.82) {$1$}; \node at (7.8,0.18) {$5e^{-4}$};
\draw[-{Latex[length=2mm]},thick] (2.25,0.25)--(4.55,0.72);
\draw[-{Latex[length=2mm]},thick] (5.35,0.72)--(7.35,0.25);
\end{tikzpicture}
\end{center}
\item La concentration est maximale au bout d'une heure.\
La concentration maximale vaut :
\[
C(1)=1.
\]
\item On calcule :
\[
C(2)=2e^{-1}=\frac2e.
\]
\item Comme $e<3$, on a $\frac1e>\frac13$.\
Donc :
\[
C(2)=\frac2e>\frac23.
\]
\item Deux heures après l'administration, la concentration est encore supérieure à $\frac23$ d'unité.
\end{enumerate}


\end{document}
